\documentclass{zhvt-classic}

% 使用本地源流明体（GenRyuMin2 TW）按文件路径加载，不依赖系统字体名
% 版框（内框／界栏／书口侧栏／鱼尾）改用与版心印章同一印泥棕色：
%   印章图片 yshs.png 实色 = RGB(135,68,52) = HTML 874434
\usepackage{xcolor}
\usepackage{tikz}
% sealbrown 双注册：xcolor 供 TikZ/pgf，L3 供 cls 版框
\definecolor{sealbrown}{HTML}{874434}
\ExplSyntaxOn
\color_set:nnn { sealbrown } { HTML } { 874434 }
\ExplSyntaxOff

\zhvtset{
  grid color       = sealbrown,
  main font        = GenRyuMin2TW-SB,
  main font path   = C:/Users/Administrator/texmf/fonts/opentype/public/GenRyuMin2TW/,
jiazhu font      = GenRyuMin2TW-M,
  jiazhu font path = C:/Users/Administrator/texmf/fonts/opentype/public/GenRyuMin2TW/,
  % 页码框容纳 4 字：正常页「九十七」+前置空 1 格=4 行恰好；
  % 100+ 页「一百零一」共 5 行超出，末字「一/二/三」被框裁掉（2026-10-04 实测），故放宽到 5
  shukou page number chars = 5,
  % 书口书名框高从默认 6 字减到 4 字，把省出的 2 字高拨给 chapter 框，
  % 使 6 字书名+卷名放进 9 字高的 chapter 框后不再压线（2026-10-05）
  shukou title chars = 4,
}

\title{史记}[绛侯周勃世家]
\author{汉·司马迁}
\maker{闲叟漫记·重製}


% ===== 索隐述赞样式（2026-11-08 四次定稿 · 愚大人亲手改定参数）=====
% ⚠ 竖排前提：自定义字体家族必须带 RawFeature={vertical:+vert}；
%   simhei.ttf 加 +vhal 会整体加载失败静默回退，只写 +vert；simfang 同理。
\newCJKfontfamily{\shuzanfont}[Path=C:/Windows/Fonts/,RawFeature={vertical:+vert}]{simhei.ttf}
\newCJKfontfamily{\shuzanbody}[Path=C:/Windows/Fonts/,RawFeature={vertical:+vert}]{simfang.ttf}
% 题头：152x30pt 圆角 6pt 棕底白字；baseline (0,2) 使整框落入界行格子不压版框线；
% 白字节点 x 不补偿、y +8.2pt（像素级墨迹实测 dx/dy≈0，L/R 留白约 20pt）。
\newcommand{\shuzanhead}[1]{%
  \begin{tikzpicture}[baseline={(0,2)}]
    \fill[sealbrown, rounded corners=6pt]
      (-76pt,-15pt) rectangle (76pt,15pt);
    \node[text=white, inner sep=0pt, anchor=center] at (0pt,8.2pt)
      {\shuzanfont\fontsize{23}{23}\selectfont #1};
  \end{tikzpicture}%
}
% ★ 修正书口卷名位置：chapter 框起始位置改回 cls 默认值，避免全局影响。
%   只在每个卷首页通过 \juan 命令里的空命令微调该页位置（安全，不会跨页污染）。
\ExplSyntaxOn
\AtBeginDocument{
  \dim_set:Nn \l__zhvt_shukou_chapter_title_start_dim
    { \l__zhvt_shukou_title_height_dim + \l__zhvt_font_size_dim * 2 }
  \dim_set:Nn \l__zhvt_shukou_page_number_start_dim
    { \l__zhvt_shukou_chapter_title_start_dim + \l__zhvt_shukou_chapter_title_height_dim }
}
\ExplSyntaxOff

\begin{document}

% ★★ 版心（书口）排布覆写：把书名放到鱼尾之下（2026-10-04）
%
%   cls 默认版心自上而下：书名（title 静态框，在鱼尾之上）→ 鱼尾 → 篇名（chapter 动态框）→ 页码
%   本模板改为：          鱼尾（位置不动）→ 书名 → 篇名 → 页码分割线 → 页码
%
%   手法：把 cls 放在 title 静态框里的书名，改放进 chapter 动态框（鱼尾之下）；
%         title 静态框留空。书名取 \l__zhvt_title_str，即 \title{...} 的第一参数，
%         所以换书时只改 \title，不必在此重复写书名。
%
%   ⚠ 三条必须遵守的写法（2026-10-04 实测）：
%     ① 汉字必须用 \hbox{...} 包住。flowfram 会把框内容装进一个宽 14.4mm（= 一个 line_sep
%        = 书口栏宽）的纵向 minipage 并按该宽度折行，而整页靠 PDF /Rotate 90 显示、
%        框内文字本身不变向 —— 不包 \hbox 时「围炉夜话」4 字会被逐字折成 4 行，
%        在版面上就成了横排 4 字并横向溢出版框。
%     ② 竖排下**一个 \hbox 就是一整列**。书名与篇名要同处一列，必须放进同一个 \hbox；
%        分成两个 \hbox 时，后一个会被当成新的一列、甩到隔壁栏去
%        （实测「序」跑到书口栏外侧 14.4mm 的正文列上）。中间的 \quad 在框内是横向空白，
%        显示出来正好是书名与篇名之间的纵向间隔。
%     ③ ★ 内容盒要写成 **与框等宽的定宽盒** `\hbox to \l__zhvt_shukou_chapter_title_height_dim{…\hfil}`。
%        实测：chapter 框是把内容盒沿纵向**居中**摆放的，内容宽度一变，书名就跟着漂：
%        自然宽度的 \hbox（书名 38.4mm）落在 x=157.2mm，加篇名后（57.6mm）落在 x=152.4mm，
%        目录页（无篇名）与正文页的书名就差半格。做成等宽盒后，无论有没有篇名，
%        内容固定从框顶起排，实测两页均落在 x=142.8mm，不再漂。
\ExplSyntaxOn
% 把书口 chapter 框（书名+篇名）起始位置下移 1 字高，避免 4 字卷名顶到鱼尾
% ⚠ cls 的 \__zhvt_calculate_dimensions: 在 \AtBeginDocument 里执行，会重算该 dim；
%   必须把覆写挂到 begin document 之后，否则会被二次覆盖（2026-10-05 实测）。
% (旧方案已移到导言区，直接修改 chapter 框起始位置)
% 卷标题命令 #1=书口篇名(鱼尾) #2=正文卷名（含卷序，卷名内空格用 ~）
%
% ★★ 目录页两项修订（2026-10-05，愚大人指出）★★
%   ① 页码起点统一对齐：
%      原状 —— 目录项写成 \par\hspace{4em}卷名\hspace{4em}页码，页码起点 = 卷名实宽 + 4em。
%      卷名实宽随字数变：「卷一~梁惠王上」6 字 +~、「卷七~离娄上」5 字 +~，
%      差 1 个字位 → 卷七~卷十 4 条的页码整体比其余 10 条高 27.21pt（一个字），看着上下不齐。
%      改法 —— 卷名补全角空格（宽 = 1 字 = 27.215pt）补足到 **7 字符位**再写进目录，
%      14 条页码起点遂全部落在同一条竖线上（起点对齐；末字随页码字数自然参差）。
%      ⚠ 只对目录里的副本 \l_tmpc_tl 补，正文卷名仍用 \l_tmpa_tl 原样 —— 正文像素不受影响。
%   ② 卷名整体上移半字位：前导 \hspace 由 4em 改 3.5em。
%      实测原状目录块在版框内偏下（天头留白 139.6pt / 地脚留白 110.7pt），
%      上移 0.5em = 13.6pt 后上下基本均分，视觉居中。
\NewDocumentCommand { \juan } { mm } {
	\clearpage
	\tl_set:Nn \l__zhvt_chapter_title_tl { #1 }
	\tl_set:Nn \l_tmpa_tl { #2 }
	\tl_set_eq:NN \l_tmpc_tl \l_tmpa_tl
	\int_step_inline:nn { 7 - \str_count:n { #2 } }
	  { \tl_put_right:Nn \l_tmpc_tl { 　 } }
	\tl_set:Nx \l_tmpb_tl { \zhnum { page } }
	\tl_put_left:Nn \l_tmpb_tl { \hspace { 4em } }
	\tl_put_left:Nn \l_tmpb_tl { \l_tmpc_tl }
	\tl_put_left:Nn \l_tmpb_tl { \par \hspace { 3.5em } }
	\addtocontents { toc } { \l_tmpb_tl }
	\skip_horizontal:n { \l__zhvt_font_size_dim * 4 }
	\l_tmpa_tl
	\par
	\thispagestyle { empty }
}
\cs_set_protected:Npn \__zhvt_set_shukou_content: {
	% ★ 鱼尾上方留白（title 静态框）嵌入闲叟钤印：
	%   整页 /Rotate 90，故此处把图用 \rotatebox{90} 转正为竖章；
	%   尺寸卡在版心条幅内（书内约 7mm 宽 × 33.8mm 高，留余量不压鱼尾）。
	\setstaticcontents * { title } {
		\vbox_to_ht:nn { \l__zhvt_baseline_skip_dim } {
			\vskip 14.2pt
			\hbox_to_wd:nn { \l__zhvt_shukou_title_height_dim } {
				\hfil
				\rotatebox [ origin = c ] { 90 } { \includegraphics [ width = 7mm , height = 33.8mm ] { yshs.png } }
				\hfil
			}
		}
	}
% ★ 版心（书口）统一改法（2026-11-16 愚大人定）：18pt＋框内居中＋书口文字同版框色
%   ① 原 27.2pt 时 7 字长篇名（如「廉颇蔺相如列传」）末字溢出 chapter 框、压住页码；
%      改 \fontsize{18pt}{18pt}，前后 \hfil 令字串在框内居中（不随篇名字数漂移）。
%   ② \vbox{\vskip 8.2pt} 定量下移，抵消缩字号后栏宽方向 8.2pt 上偏（像素实测 y 中心 543.3 = 原版）。
%      ⚠ \raisebox / \vspace* 在 flowfram 框内无效；换字号须同步换补偿：16pt→10.1、19pt→7.4。
%   ③ chapter 与 page~number 均置 \color{sealbrown}（HTML 874434），与版框／界行同色；正文仍黑。
%   ⚠ 页码字位、钤印、鱼尾、界行、正文一律未动；五部合订本同此口径。
	\setdynamiccontents * { chapter }
	{ \color{sealbrown} \vbox { \vskip 8.2pt \hbox to \l__zhvt_shukou_chapter_title_height_dim
		{ \hfil \fontsize{18pt}{18pt}\selectfont \l__zhvt_title_str \quad \l__zhvt_chapter_title_tl \hfil } } }
	\setdynamiccontents * { page~number } { \color{sealbrown} \quad \zhnum { page } }
}
\__zhvt_set_shukou_content:
\ExplSyntaxOff

\maketitle

% 单篇，无多卷目录，不渲染目录页

\juan{绛侯周勃世家}{卷五十七~绛侯周勃世家第二十七}
\hspace*{2em}绛侯周勃者。沛人也。其先卷人。徙沛。勃以织薄曲为生。常为人吹箫给丧事。材官引彊。\\
\hspace*{2em}高祖之为沛公初起。勃以中涓从攻胡陵。下方与。方与反。与战。卻適。攻丰。击秦军砀东。还军留及萧。复攻砀。破之。下下邑。先登。赐爵五大夫。攻蒙。虞。取之。击章邯车骑。殿。定魏地。攻爰戚。东缗。以往至栗。取之。攻齧桑。先登。击秦军阿下。破之。追至濮阳。下甄城。攻都关。定陶。袭取宛朐。得单父令。夜袭取临济。攻张。以前至卷。破之。击李由军雍丘下。攻开封。先至城下为多。后章邯破杀项梁。沛公与项羽引兵东如砀。自初起沛还至砀。一岁二月。楚怀王封沛公号安武侯。为砀郡长。沛公拜勃为虎贲令。以令从沛公定魏地。攻东郡尉於城武。破之。击王离军。破之。攻长社。先登。攻颍阳。缑氏。绝河津。击赵贲军尸北。南攻南阳守齮。破武关。峣关。破秦军於蓝田。至咸阳。灭秦。\\
\hspace*{2em}项羽至。以沛公为汉王。汉王赐勃爵为威武侯。从入汉中。拜为将军。还定三秦。至秦。赐食邑怀德。攻槐里。好畤。最。击赵贲。内史保於咸阳。最。北攻漆。击章平。姚卬军。西定汧。还下郿。频阳。围章邯废丘。破西丞。击盗巴军。破之。攻上邽。东守峣关。转击项籍。攻曲逆。最。还守敖仓。追项籍。籍已死。因东定楚地泗。东海郡。凡得二十二县。还守雒阳。栎阳。赐与颍侯共食锺离。以将军从高帝反者燕王臧荼。破之易下。所将卒当驰道为多。赐爵列侯。剖符世世勿绝。食绛八千一百八十户。号绛侯。\\
\hspace*{2em}以将军从高帝击反韩王信於代。降下霍人。以前至武泉。击胡骑。破之武泉北。转攻韩信军铜鞮。破之。还。降太原六城。击韩信胡骑晋阳下。破之。下晋阳。后击韩信军於硰石。破之。追北八十里。还攻楼烦三城。因击胡骑平城下。所将卒当驰道为多。勃迁为太尉。\\
\hspace*{2em}击陈豨。屠马邑。所将卒斩豨将军乘马絺。击韩信。陈豨。赵利军於楼烦。破之。得豨将宋最。雁门守。因转攻得云中守遫。丞相箕肆。将勋。定雁门郡十七县。云中郡十二县。因复击豨灵丘。破之。斩豨。得豨丞相程纵。将军陈武。都尉高肆。定代郡九县。燕王卢绾反。勃以相国代樊哙将。击下蓟。得绾大将抵。丞相偃。守陉。太尉弱。御史大夫施。屠浑都。破绾军上兰。复击破绾军沮阳。追至长城。定上谷十二县。右北平十六县。辽西。辽东二十九县。渔阳二十二县。最从高帝得相国一人。丞相二人。将军。二千石各三人。别破军二。下城三。定郡五。县七十九。得丞相。大将各一人。\\
\hspace*{2em}勃为人木彊敦厚。高帝以为可属大事。勃不好文学。每召诸生说士。东乡坐而责之。﹁趣为我语。﹂其椎少文如此。\\
\hspace*{2em}勃既定燕而归。高祖已崩矣。以列侯事孝惠帝。孝惠帝六年。置太尉官。以勃为太尉。十岁。高后崩。吕禄以赵王为汉上将军。吕产以吕王为汉相国。秉汉权。欲危刘氏。勃为太尉。不得入军门。陈平为丞相。不得任事。於是勃与平谋。卒诛诸吕而立孝文皇帝。其语在吕后。孝文事中。\\
\hspace*{2em}文帝既立。以勃为右丞相。赐金五千斤。食邑万户。居月馀。人或说勃曰。﹁君既诛诸吕。立代王。威震天下。而君受厚赏。处尊位。以宠。久之即祸及身矣。﹂勃惧。亦自危。乃谢请归相印。上许之。岁馀。丞相平卒。上复以勃为丞相。十馀月。上曰。﹁前日吾诏列侯就国。或未能行。丞相吾所重。其率先之。﹂乃免相就国。\\
\hspace*{2em}岁馀。每河东守尉行县至绛。绛侯勃自畏恐诛。常被甲。令家人持兵以见之。其后人有上书告勃欲反。下廷尉。廷尉下其事长安。逮捕勃治之。勃恐。不知置辞。吏稍侵辱之。勃以千金与狱吏。狱吏乃书牍背示之。曰﹁以公主为证﹂。公主者。孝文帝女也。勃太子胜之尚之。故狱吏教引为证。勃之益封受赐。尽以予薄昭。及系急。薄昭为言薄太后。太后亦以为无反事。文帝朝。太后以冒絮提文帝。曰。﹁绛侯绾皇帝玺。将兵於北军。不以此时反。今居一小县。顾欲反邪。﹂文帝既见绛侯狱辞。乃谢曰。﹁吏方验而出之。﹂於是使使持节赦绛侯。复爵邑。绛侯既出。曰。﹁吾尝将百万军。然安知狱吏之贵乎。﹂\\
\hspace*{2em}绛侯复就国。孝文帝十一年卒。谥为武侯。子胜之代侯。六岁。尚公主。不相中。坐杀人。国除。绝一岁。文帝乃择绛侯勃子贤者河内守亚夫。封为条侯。续绛侯后。\\
\hspace*{2em}条侯亚夫自未侯为河内守时。许负相之。曰。﹁君后三岁而侯。侯八岁为将相。持国秉。贵重矣。於人臣无两。其后九岁而君饿死。﹂亚夫笑曰。﹁臣之兄已代父侯矣。有如卒。子当代。亚夫何说侯乎。然既已贵如负言。又何说饿死。指示我。﹂许负指其口曰。﹁有从理入口。此饿死法也。﹂居三岁。其兄绛侯胜之有罪。孝文帝择绛侯子贤者。皆推亚夫。乃封亚夫为条侯。续绛侯后。\\
\hspace*{2em}文帝之后六年。匈奴大入边。乃以宗正刘礼为将军。军霸上。祝兹侯徐厉为将军。军棘门。以河内守亚夫为将军。军细柳。以备胡。上自劳军。至霸上及棘门军。直驰入。将以下骑送迎。已而之细柳军。军士吏被甲。锐兵刃。彀弓弩。持满。天子先驱至。不得入。先驱曰。﹁天子且至。﹂军门都尉曰。﹁将军令曰﹃军中闻将军令。不闻天子之诏﹄。﹂居无何。上至。又不得入。於是上乃使使持节诏将军。﹁吾欲入劳军。﹂亚夫乃传言开壁门。壁门士吏谓从属车骑曰。﹁将军约。军中不得驱驰。﹂於是天子乃按辔徐行。至营。将军亚夫持兵揖曰。﹁介胄之士不拜。请以军礼见。﹂天子为动。改容式车。使人称谢。﹁皇帝敬劳将军。﹂成礼而去。既出军门。群臣皆惊。文帝曰。﹁嗟乎。此真将军矣。曩者霸上。棘门军。若兒戏耳。其将固可袭而虏也。至於亚夫。可得而犯邪。﹂称善者久之。月馀。三军皆罢。乃拜亚夫为中尉。\\
\hspace*{2em}孝文且崩时。诫太子曰。﹁即有缓急。周亚夫真可任将兵。﹂文帝崩。拜亚夫为车骑将军。\\
\hspace*{2em}孝景三年。吴楚反。亚夫以中尉为太尉。东击吴楚。因自请上曰。﹁楚兵剽轻。难与争锋。原以梁委之。绝其粮道。乃可制。﹂上许之。\\
\hspace*{2em}太尉既会兵荥阳。吴方攻梁。梁急。请救。太尉引兵东北走昌邑。深壁而守。梁日使使请太尉。太尉守便宜。不肯往。梁上书言景帝。景帝使使诏救梁。太尉不奉诏。坚壁不出。而使轻骑兵弓高侯等绝吴楚兵后食道。吴兵乏粮。饥。数欲挑战。终不出。夜。军中惊。内相攻击扰乱。至於太尉帐下。太尉终卧不起。顷之。复定。后吴奔壁东南陬。太尉使备西北。已而其精兵果奔西北。不得入。吴兵既饿。乃引而去。太尉出精兵追击。大破之。吴王濞弃其军。而与壮士数千人亡走。保於江南丹徒。汉兵因乘胜。遂尽虏之。降其兵。购吴王千金。月馀。越人斩吴王头以告。凡相攻守三月。而吴楚破平。於是诸将乃以太尉计谋为是。由此梁孝王与太尉有卻。\\
\hspace*{2em}归。复置太尉官。五岁。迁为丞相。景帝甚重之。景帝废栗太子。丞相固争之。不得。景帝由此疏之。而梁孝王每朝。常与太后言条侯之短。\\
\hspace*{2em}窦太后曰。﹁皇后兄王信可侯也。﹂景帝让曰。﹁始南皮。章武侯先帝不侯。及臣即位乃侯之。信未得封也。﹂窦太后曰。﹁人主各以时行耳。自窦长君在时。竟不得侯。死后乃其子彭祖顾得侯。吾甚恨之。帝趣侯信也。﹂景帝曰。﹁请得与丞相议之。﹂丞相议之。亚夫曰。﹁高皇帝约﹃非刘氏不得王。非有功不得侯。不如约。天下共击之﹄。今信虽皇后兄。无功。侯之。非约也。﹂景帝默然而止。\\
\hspace*{2em}其后匈奴王徐卢等五人降。景帝欲侯之以劝后。丞相亚夫曰。﹁彼背其主降陛下。陛下侯之。则何以责人臣不守节者乎。﹂景帝曰。﹁丞相议不可用。﹂乃悉封徐卢等为列侯。亚夫因谢病。景帝中三年。以病免相。\\
\hspace*{2em}顷之。景帝居禁中。召条侯。赐食。独置大胾。无切肉。又不置櫡。条侯心不平。顾谓尚席取櫡。景帝视而笑曰。﹁此不足君所乎。﹂条侯免冠谢。上起。条侯因趋出。景帝以目送之。曰。﹁此怏怏者非少主臣也。﹂\\
\hspace*{2em}居无何。条侯子为父买工官尚方甲楯五百被可以葬者。取庸苦之。不予钱。庸知其盗买县官器。怒而上变告子。事连汙条侯。书既闻上。上下吏。吏簿责条侯。条侯不对。景帝骂之曰。﹁吾不用也。﹂召诣廷尉。廷尉责曰。﹁君侯欲反邪。﹂亚夫曰。﹁臣所买器。乃葬器也。何谓反邪。﹂吏曰。﹁君侯纵不反地上。即欲反地下耳。﹂吏侵之益急。初。吏捕条侯。条侯欲自杀。夫人止之。以故不得死。遂入廷尉。因不食五日。呕血而死。国除。\\
\hspace*{2em}绝一岁。景帝乃更封绛侯勃他子坚为平曲侯。续绛侯后。十九年卒。谥为共侯。子建德代侯。十三年。为太子太傅。坐酎金不善。元鼎五年。有罪。国除。\\
\hspace*{2em}条侯果饿死。死后。景帝乃封王信为盖侯。\\
\hspace*{2em}太史公曰。绛侯周勃始为布衣时。鄙朴人也。才能不过凡庸。及从高祖定天下。在将相位。诸吕欲作乱。勃匡国家难。复之乎正。虽伊尹。周公。何以加哉。亚夫之用兵。持威重。执坚刃。穰苴曷有加焉。足己而不学。守节不逊。终以穷困。悲夫。\\

% ===== 索隐述赞（司马贞《史记索隐》）· 统一样式 =====
% 题头「索隐述赞曰」五字并入反白框；正文每四句一列，定宽盒 \hfil 居列中；
% \fontsize{24}{41} 列距对齐全书界行（121+41k）；\clearpage 使述赞独占一页。
\clearpage
\shuzanhead{索隐述赞曰}\linebreak
{\shuzanbody\fontsize{24}{41}\selectfont
\hbox to \hsize{\hfil 绛侯佐汉。质厚敦笃。始击砀东。亦围尸北。\hfil}\linebreak
\hbox to \hsize{\hfil 所攻必取。所讨咸克。陈豨伏诛。臧荼破国。\hfil}\linebreak
\hbox to \hsize{\hfil 事居送往。推功伏德。列侯还第。太尉下狱。\hfil}\linebreak
\hbox to \hsize{\hfil 继相条侯。绍封平曲。惜哉贤将。父子代辱。\hfil}%
\par}
\end{document}
